\subsection{问题背景}

题目研究有限时频资源下的装备用频编排。系统提供 100 个离散频段，现有 150 条用频计划分为 A、B、C 三类，优先级依次降低。每条计划由频段区间、首次使用时间区间、间隔时长和使用次数定义，并沿时间轴重复占用同一频段；若两条计划的频段区间相交，且至少有一对重复事件的时间区间也相交，则发生时频冲突。因此，判定冲突不能只比较首次时间区间，而要先还原完整重复事件，再在统一的时间—频率边界内处理计划之间的资源关系。

\subsection{问题提出}

前三问构成一条依赖链：先识别冲突，再按有限调整规则和类别优先级消解冲突，最后固定方案并计算新增 C 类计划。各问的“最多”均限定在题目给定的资源域、动作模板和冻结输入内；第四问允许改变 C 类计划的间隔时长，本阶段暂不运行。

\CUMCMQuestionLead{一}
对原始计划逐条展开其全部重复使用事件，并判定任意两条计划是否存在冲突。冲突采用两个维度同时满足的口径：频段区间有交集，且至少有一对重复事件的时间区间有交集。时间区间使用半开形式，左端点计入、右端点不计入，以避免首尾相接被误判为重叠。该问的直接输出是计划层的冲突装备对清单；事件级重叠次数作为辅助核验量保留。除计划对枚举外，还需要用离散时频资源格重新检查同一结果，从而区分“冲突装备对数”和“发生重叠的事件次数”这两个统计口径。

\CUMCMQuestionLead{二}
在问题一得到的冲突背景下，为每条原计划选择一种动作：保持不变、频率平移、时间平移或撤销。频率平移量和时间平移量分别受到给定范围限制，一条计划不能同时改变两个维度，调整后全部重复事件仍须落在资源域内，并且所有计划之间不得共享同一时频资源格。题面只说明“尽量少撤销、尽量少调整并优先保护高等级装备”，没有给出可直接相加的数值权重，因此需要把这些要求转化为可审计的条件字典序。输出不仅包括一份无冲突动作表，还要给出每一层优先目标的可行解、下界和不可行阈值证据；其中末级平移量若只得到当前可行上界，必须单独标注其证明边界。

\CUMCMQuestionLead{三}
固定问题二通过验证的具体方案，不再重新移动或撤销其中的既有计划。在同一时间—频率资源域内，按照附件给出的 C 类装备共同结构生成所有可能的新增起点，逐个展开其完整重复事件，删除与固定问题二背景发生冲突的候选，再从剩余候选中选出数量最多且彼此不冲突的一组。该问的输出是最大新增数量及相应新增计划表；“最大”仅针对当前冻结的 Q2 方案、当前资源域和 C 类模板成立。为避免结果依赖某一种约束编码，需要用资源格集合装填模型和候选两两冲突模型交叉验证最大值，并单独检查新增计划与原有方案之间的边界和重复占用。

\CUMCMQuestionLead{四}
第四问允许改变 C 类计划的间隔时长，会重新改变事件递推关系；本阶段暂不运行，因此不改写问题一至问题三的输入和结论。
